% Results Explanation for Figure 1: Model Preference Heterogeneity
% For Methodology/Results Section

\subsection{Model Preference Heterogeneity}
\label{sec:semantic_structure}

Before presenting our routing algorithm, we establish the empirical motivation for contextual routing: model preference varies by prompt, and the router's features predict this variation.

\paragraph{Why this experiment matters.}
BanditGPT routes prompts to models using contextual features extracted by \texttt{FeatureService}: sentence embeddings compressed to 32 dimensions via PCA, plus a bias term (33 dimensions total).
LinUCB then uses these features to predict per-model reward and select the best model for each prompt.
If the PCA features contain no signal about model preference, the bandit reduces to random assignment---contextual routing would be unnecessary, and a static policy (always pick one model) would suffice.
This experiment tests whether the features predict model preference on held-out data, and whether this prediction exceeds what a random projection achieves.

\paragraph{Setup.}
We evaluate on a two-model topology (Mixtral-8x7b-instruct vs.\ GPT-4-Turbo) that matches the RouteLLM benchmark~\cite{ong2024routellm}.
We project $N=750$ held-out prompts through the production PCA (trained on 80K RouteLLM battle prompts---an independent dataset) and compute the Spearman rank correlation between PC1 and the reward gap $(R_{\text{GPT-4-Turbo}} - R_{\text{Mixtral}})$.

\noindent\textbf{Data independence.}
The PCA training data (RouteLLM battles) and the held-out evaluation data (LMSYS general prompts) originate from \emph{different data sources, different sampling periods, and different prompt populations}.
The PCA has never seen any evaluation prompt.

\paragraph{Primary metric.}
We use Spearman rank correlation ($\rho$) as the primary metric because it directly tests the question the bandit needs answered: do the features predict the direction of model preference?
This matches how BanditGPT actually works---LinUCB performs regression on features to predict reward---rather than testing an artificial clustering that the router never uses.

\paragraph{Null baseline.}
To rule out projection artifacts (any 384D-to-2D projection might find \emph{some} correlated direction), we compute $|\rho|$ for 100 independent random orthonormal projections ($384 \to 2$, QR-decomposed) and compare the Router PCA $|\rho|$ to this null distribution.

\paragraph{Results.}
The router's PC1 shows a significant rank correlation with model preference:
\begin{itemize}[noitemsep,topsep=2pt]
    \item \textbf{Spearman $\rho$}: reported in Figure~\ref{fig:lmsys_holdout_structure} ($p < 0.0001$, $N=750$).
    \item \textbf{Null comparison}: the Router PCA $|\rho|$ exceeds all 100 random projections, with a signal ratio reported in the figure.
\end{itemize}

\noindent
The correlation is negative: higher PC1 values correspond to prompts where Mixtral outperforms GPT-4-Turbo.
This means the PCA's first component separates prompts by model preference direction---exactly the signal a contextual bandit needs to learn routing.

\paragraph{Model preference is heterogeneous.}
The reward gap distribution confirms that model preference is not uniform:
the majority of prompts are ties (neither model is clearly better), while a minority show strong preference for one model.
If preference were uniform, a static policy would be optimal; the heterogeneity is what makes contextual routing valuable.

\paragraph{Practical implications.}
\begin{enumerate}[noitemsep,topsep=2pt]
    \item \textbf{Contextual routing is justified.} The features predict model preference significantly better than chance. A bandit using these features can learn to route.

    \item \textbf{Most prompts are ties.} The majority of prompts show no strong model preference, meaning the cheaper model can serve them with comparable quality. This is the primary source of cost savings.

    \item \textbf{The signal is a lower bound.} PC1 is one dimension; the router uses all 32 PCA components. The full feature vector likely captures more predictive signal than PC1 alone.

    \item \textbf{No contamination.} The PCA was trained on RouteLLM battles, an entirely separate dataset from the held-out evaluation prompts.

    \item \textbf{Effect sizes are holdout-specific.} The holdout was constructed via stratified sampling including a difficulty dimension derived from oracle rewards, which enriches it for prompts with clear model preferences (Appendix~\ref{app:stratification}). The \emph{existence} of the signal is robust; the \emph{magnitude} in deployment will depend on the user's prompt distribution.
\end{enumerate}

\paragraph{Methodological notes.}
\begin{itemize}[noitemsep,topsep=2pt]
    \item \textbf{Data:} Held-out test set only ($N=750$). Dev set excluded.
    \item \textbf{Rewards:} Discrete pairwise outcomes from LMSYS Chatbot Arena.
    \item \textbf{Metric:} Spearman $\rho$ (rank correlation), which is appropriate for ordinal/discrete data and makes no distributional assumptions.
    \item \textbf{Null baseline:} 100 random orthonormal projections (QR-decomposed), providing the null distribution of $|\rho|$ under the hypothesis that the PCA direction is no more informative than chance.
    \item \textbf{Prompt-type mechanism:} The correlation reflects prompt-type variation that correlates with model preference (e.g., instruction-following templates favor Mixtral). The same PC1 axis may not separate winners identically for a different model pair.
\end{itemize}

\paragraph{Connection to routing evaluation.}
This experiment establishes a \emph{necessary condition}: the router's features predict model preference.
Whether the bandit \emph{exploits} this signal for practical routing gains is tested in the routing evaluation (Table~\ref{tab:results}).
